\chapter{Theoretical Part}

\section{Overview of Smart Plant Monitoring Systems}

This section should explain what smart plant monitoring systems are, how they differ from manual observation, and why they are useful in precision agriculture, indoor gardening, and automated environmental control. You can define the concept of cyber-physical monitoring systems and explain how data-driven care decisions improve stability and reduce human error.

It is useful to include a short literature-based overview of system types, for example single-sensor hobby systems, cloud-connected commercial systems, and research prototypes integrating decision-support logic \cite{wolfert2017smartfarming,elijah2018iot}.

\section{IoT Technologies in Plant Condition Monitoring}

This section should describe the role of the Internet of Things in environmental sensing, data acquisition, and automated monitoring. Explain that IoT enables devices such as ESP32-based controllers to collect sensor values and send telemetry to a remote application or dashboard in real time.

\subsection{Sensors for Environmental Data Collection}

In this subsection, describe the physical parameters important for plant condition assessment. For the implemented system, these include soil moisture, air temperature, humidity, and light level. Explain why each parameter matters biologically and operationally.

You can also mention common low-cost sensor categories, measurement limitations, calibration concerns, and sources of noise in real environments \cite{jawad2017wsn,ayaz2019iotagriculture}.

\subsection{Data Transmission and Real-Time Monitoring}

This subsection should explain how sensor devices communicate with the backend. Describe the idea of periodic telemetry transmission over HTTP, secure device authentication, and real-time or near-real-time data visibility in the user dashboard.

You may describe typical communication flow as follows: sensor reading $\rightarrow$ network transmission $\rightarrow$ backend API $\rightarrow$ database storage $\rightarrow$ dashboard visualization.

\section{Artificial Intelligence for Plant Condition Analysis}

This section should explain the role of artificial intelligence and decision-support methods in turning raw environmental values into useful plant-care conclusions. Emphasize that AI in this project is applied to condition classification and recommendation support rather than image-based disease diagnosis.

\subsection{Rule-Based Analysis and Threshold Logic}

In this subsection, explain the theoretical basis of expert-defined thresholds and explainable condition scoring. Describe how a system can classify plant condition by comparing measured values against acceptable ranges and by identifying deviations and risk factors.

You can connect this subsection directly to the implemented baseline logic in the final system: condition scoring, health score, risk factors, confidence estimation, and human-readable explanation.

\subsection{Machine Learning for Plant Condition Classification}

This subsection should introduce supervised machine learning for multi-class classification. Explain the idea of training a model on labeled environmental data to predict classes such as \texttt{normal}, \texttt{attention}, and \texttt{critical}. Since the final prototype uses a Random Forest classifier, this is the place to introduce ensemble trees, feature importance, robustness, and explainability \cite{breiman2001randomforests,liakos2018machinelearning}.

You can also explain why Random Forest is suitable for diploma projects:
\begin{itemize}
    \item it works well on structured tabular data;
    \item it is easier to explain than deep neural models;
    \item it is stable on small and medium datasets;
    \item it provides probabilistic class estimates for decision support.
\end{itemize}

\subsection{AI-Based Recommendations for Plant Care}

This subsection should describe how a system converts analysis results into actionable care advice. Explain that recommendations may be generated from explicit risk factors, model-supported condition class, or hybrid logic. In the implemented system, recommendations remain rule-based for safety and explainability, while machine learning adds supporting classification evidence.

\section{Review of Existing Systems and Their Limitations}

This section should compare several existing plant monitoring or smart agriculture systems from the literature and explain their limitations. Typical limitations may include lack of explainability, weak integration between sensing and user dashboards, dependence on cloud services, poor security, or absence of intelligent recommendations \cite{mishra2021plantmonitoring,elijah2018iot}.

Suggested comparison criteria:
\begin{itemize}
    \item monitored parameters;
    \item connectivity model;
    \item presence or absence of mobile/web dashboard;
    \item support for alerts;
    \item use of AI or ML;
    \item security model;
    \item scalability and cost.
\end{itemize}

\section{Requirements for the Proposed System}

This section should translate the theoretical review into concrete functional and non-functional requirements for the proposed diploma system. Explain what the final prototype must do and what quality constraints it must satisfy.

Examples of functional requirements:
\begin{itemize}
    \item user registration and authentication;
    \item plant profile creation and management;
    \item sensor registration and secure ingest;
    \item dashboard visualization of current and historical readings;
    \item condition analysis, recommendation generation, and alerts.
\end{itemize}

Examples of non-functional requirements:
\begin{itemize}
    \item reliability of data ingestion;
    \item clarity and explainability of analysis;
    \item basic security for user and device access;
    \item maintainability and modularity of the software architecture.
\end{itemize}
